\documentclass[11pt,aps,preprint,nofootinbib,superscriptaddress,tightenlines]{revtex4-2}
\usepackage[T1]{fontenc}
\usepackage{amsmath,amssymb,mathtools}
\let\lambdabar\relax   % revtex defines it; newpxmath does too
\usepackage{newpxtext,newpxmath}
\usepackage{bm}
\usepackage[final]{microtype}
\usepackage{xcolor}
\usepackage{booktabs}
\usepackage{graphicx}
\usepackage{dsfont}
\usepackage[most]{tcolorbox}
\newcommand{\one}{\mathds{1}}
\newcommand{\avg}[1]{\langle #1\rangle}
\definecolor{darkblue}{RGB}{20,60,140}
\definecolor{accent}{RGB}{30,80,150}
\definecolor{tint}{RGB}{243,246,251}
\newtcolorbox{keyeq}{enhanced, colback=tint, frame hidden, borderline west={2pt}{0pt}{accent},
  sharp corners, left=4pt, right=4pt, top=0pt, bottom=0pt, before skip=6pt, after skip=6pt}
\usepackage[colorlinks=true,linkcolor=darkblue,citecolor=darkblue,urlcolor=darkblue]{hyperref}

\def\st{\begin{equation}}
\def\stp{\end{equation}}
\def\Eq#1{eq.~(\ref{#1})}
\def\Sect#1{Sect.~\ref{#1}}
\def\Fig#1{Fig.~\ref{#1}}
\def\Tab#1{Table~\ref{#1}}
\def\App#1{Appendix~\ref{#1}}
\def\tr{\operatorname{tr}}
\let\oldparagraph\paragraph
\renewcommand{\paragraph}{\par\bigskip\oldparagraph}



\begin{document}
\sloppy

\title{The state-counting renormalization of the particle model}
\author{D.~Teaney}
\affiliation{Department of Physics and Astronomy, Stony Brook University,
Stony Brook, New York 11794, USA}
\date{\today}

\begin{abstract}
This note summarizes the renormalization scheme that we plan to implement.  It applies to any equation of state $p(T_{\rm rel})$ and to neutral
labels ($\mathcal B=0$).  The entropy of a particle is defined as the logarithm of the number of states below its energy, and the number of states
is fixed by requiring that one particle at fixed volume has the partition function of the target fluid.  At large volume the number of states is
given by a saddle point.  The first term of the saddle-point expansion defines \emph{state counting at one loop}.  In this scheme the entropy of a
particle is the target entropy less two counterterms that are evaluated at the target's own state.  One comes from the determinant of the
susceptibility and the other from the labels, and both are known in closed form.  The temperature,
the velocity and the pressure of a particle are then those of the target state with small shifts, and only the two particle-level maps of the code
change.  The next term of the expansion, of order $1/V$, is also known in closed form in terms of $c_s^2$ and its first two derivatives.  We record
it in an appendix but do not plan to implement it soon.  The longer derivations, the exact series and the comparison with the other schemes are in
\nolinkurl{renorm2.tex}.  Written 2026-11-24.
\end{abstract}

\maketitle
\tableofcontents

%=============================================================================
\section{The condition}
\label{sec:condition}

A particle has a label $s$, a momentum $\bm m$ and a volume $V$, the Liouville measure $e^{s}ds\,d^dm$, and a Hamiltonian $H(s,\bm m,V)$.  At
fixed positions the particles are independent (\nolinkurl{partition.tex}), so the model is fixed by the partition function of one particle.  We
require that it be the target's, with the factor that cancels the ideal gas of labels,
\begin{keyeq}
\st
\int ds\,d^dm\;e^{\,s-\beta^0H(s,\bm m,V)+\bm b\cdot\bm m}=\frac{e^{\,p\,\beta^0V}}{e\,V}
\label{condition}
\stp
\end{keyeq}
for every $\beta^0=1/T$ and $\bm b=\bm v/T$.  When every particle has the mean volume $\bar V$, the partition function of the model is then
\st
\frac{\ln\mathcal Z_N}N=\ln\bar V+1+\ln Z_1=p\,\beta^0\bar V .
\stp
The model therefore has the pressure, the energy and the momentum of the target, and the
chemical potential of the labels is zero.  The fluctuations of the volumes correct this at order $1/\bar nh^2$, and that correction is not
addressed in this note.

\paragraph{The number of states.}  We change the variable of integration from the label to the energy, $e=H(s,\bm m,V)$, at fixed $\bm m$ and $V$.
Because the measure is $e^sds=d(e^s)$, the quantity $e^{s}$ is the number of states of the particle below the energy $e$.  In the new variable
\Eq{condition} reads
\st
\int de\,d^dm\;\frac{\partial e^{s}}{\partial e}\;e^{-\beta^0e+\bm b\cdot\bm m}=\frac{e^{\,p\beta^0V}}{eV} ,
\stp
so the density of states $g$ is the inverse Laplace transform of the target, and the label is the logarithm of its integral over the energy,
\begin{keyeq}
\st
e^{\,p\,\beta^0V}=\int de\,d^dm\;g(e,\bm m,V)\,e^{-\beta^0e+\bm b\cdot\bm m},\qquad
G(e,\bm m,V)=\int^{e}\!g\,de' ,
\stp
\st
s=S(e,\bm m,V)\equiv\ln G-\ln V-1 .
\label{Sdef}
\stp
\end{keyeq}
We call $S(e,\bm m,V)$ the entropy function of the particle.  The Hamiltonian is obtained by solving $s=S(e,\bm m,V)$ for the energy.  There is no
freedom in this construction, since the pressure $p(\beta)$ is the input and the function $S$ is the output.  The derivatives of $S$ are the first
law of the particle,
\begin{keyeq}
\st
dS=\frac{de-\bm v_a\cdot d\bm m+P_a\,dV}{T_a},\qquad
T_a=\frac{\partial H}{\partial s},\qquad \bm v_a=\frac{\partial H}{\partial\bm m}=\dot{\bm x}_a,\qquad P_a=-\frac{\partial H}{\partial V} ,
\label{firstlaw}
\stp
\end{keyeq}
so the temperature, the velocity and the pressure of a particle follow from $S$ by differentiation.  In terms of the number of states and its
density, $g=\partial G/\partial e$, they are
\begin{keyeq}
\st
T_a=\frac Gg,\qquad \bm v_a=-\frac1g\,\frac{\partial G}{\partial\bm m},\qquad P_a=\frac1g\,\frac{\partial G}{\partial V}-\frac{T_a}V .
\label{fromG}
\stp
\end{keyeq}
The temperature of a particle is the number of states below its energy divided by the density of states at its energy.  The last term of the
pressure comes from the factor $1/V$ in \Eq{condition}.

\paragraph{The weight in the energy.}  When the energy is used as the variable in place of the label, the measure becomes
\st
e^{s}\,ds=\frac{\partial e^{S}}{\partial e}\,de=\frac{e^{S}}{T_a}\,de .
\stp
The weight of a particle in the energy is therefore $S-\ln T_a$, which is $\ln g$ up to the terms in
the volume.  This is the weight that the viscous step uses (\nolinkurl{scheme.tex}).

%=============================================================================
\section{The target}
\label{sec:target}

\paragraph{Notation.}  $P^\mu=(e,\bm m)$ and $\beta_\mu=(-\beta^0,\bm b)$, so that $\beta_\mu P^\mu=-\beta^0e+\bm b\cdot\bm m$; the metric is mostly
plus, $-\beta_\mu\beta^\mu=1/T_{\rm rel}^2$, $T=T_{\rm rel}/\gamma$ is the bare temperature, and $D=d+1$.  The equation of state enters through
\st
q\equiv\frac1{c_s^2}=\frac{d\ln s}{d\ln T_{\rm rel}},\qquad \dot q\equiv\frac{dq}{d\ln T_{\rm rel}},\qquad \ddot q ,
\stp
with $w=T_{\rm rel}s$ and $\varepsilon=w-p$ the rest-frame enthalpy and energy densities.

\paragraph{One scalar.}  $p\,\beta^\mu$ is the gradient of a function $\Psi$ of $\beta_\mu\beta^\mu$, and the target is
\st
\ln Z_{\rm t}=V\,\Phi,\qquad \Phi\equiv p\,\beta^0=\frac{\partial\Psi}{\partial\beta_0},\qquad
T^{\mu\nu}=\frac{\partial^2\Psi}{\partial\beta_\mu\partial\beta_\nu},\qquad
\chi^{\mu\nu}\equiv\frac{\partial^2\Phi}{\partial\beta_\mu\partial\beta_\nu}=\frac{\partial T^{0\mu}}{\partial\beta_\nu} .
\stp
$\chi^{\mu\nu}$ is the matrix of susceptibilities of a cell, $\avg{\delta P^\mu\delta P^\nu}=V\chi^{\mu\nu}$, and
\st
\det\chi^{\mu\nu}=\gamma^{D}\,(wT_{\rm rel})^{D}\,a,\qquad a\equiv q-v^2=\frac{1-c_s^2v^2}{c_s^2} .
\label{det}
\stp
The quantity $a$ is the acoustic norm of the time normal $n_\mu$ of the cell.  In terms of the acoustic metric it is
\st
\gamma^2a=-G_{\rm ac}^{\mu\nu}n_\mu n_\nu,\qquad G_{\rm ac}^{\mu\nu}=g^{\mu\nu}+\Big(1-\frac1{c_s^2}\Big)u^\mu u^\nu .
\stp
It lies between $q-1$ and $q$ for $v<1$, so
it never vanishes.

\paragraph{The target state of a particle.}  The target is the fluid that the model should reproduce.  It is described by the input equation of state, without any
correction from the fluctuations of a particle.  Given the energy, the momentum and the volume of a particle, its \emph{target state} is the
temperature $T_{\rm rel}$ and the velocity $\bm v$ that this fluid would have if a cell of volume $V$ held exactly the energy $e$ and the momentum
$\bm m$.  This is the condition $P^\mu=V\,T^{0\mu}$, or
\st
e=V\big(w\gamma^2-p\big),\qquad \bm m=V\,w\gamma^2\bm v,\qquad S_{\rm t}(e,\bm m,V)=\gamma\,s(T_{\rm rel})\,V .
\label{targetstate}
\stp
The velocity is along $\bm m$.  With the rapidity $\theta$, $v=\tanh\theta$ and $\gamma=\cosh\theta$, the thermal vector of the target state is
\st
\bar\beta_\mu=\frac{(-\cosh\theta,\ \sinh\theta\,\hat{\bm m})}{T_{\rm rel}} .
\stp
This is the state the code assigns to a particle at tree level, by the
inversion \texttt{eos\_from\_energy}.  $S_{\rm t}$ is the entropy of the target fluid in the cell, and its differential is the first law of the target,
\st
dS_{\rm t}=-\bar\beta_\mu\,dP^\mu+p\,\bar\beta^0\,dV .
\stp

%=============================================================================
\section{State counting at one loop}
\label{sec:oneloop}

\paragraph{The saddle point.}  By \Eq{Sdef}, $e^{V\Phi}$ is the Laplace transform of the density of states, so $g$ is the inverse transform,
\st
g(e,\bm m,V)=\int\frac{d^D\beta}{(2\pi i)^D}\;\exp\big[V\Phi(\beta)-\beta_\mu P^\mu\big],\qquad -\beta_\mu P^\mu=\beta^0e-\bm b\cdot\bm m .
\stp
The number of states is its integral over the energy at fixed momentum.  The energy appears only in the factor $e^{\beta^0e}$, whose integral is
\st
\int^{e}e^{\beta^0e'}\,de'=\frac{e^{\beta^0e}}{\beta^0} .
\stp
The lower limit does not contribute because the contour has ${\rm Re}\,\beta^0>0$.  The integration
over the energy therefore puts one factor of $1/\beta^0$ into the integrand,
\st
G(e,\bm m,V)=\int^{e}\!g\,de'=\int\frac{d^D\beta}{(2\pi i)^D}\;\frac1{\beta^0}\,\exp\big[V\Phi(\beta)-\beta_\mu P^\mu\big] .
\stp
The factor $1/(eV)$ of \Eq{condition} does not depend on $\beta$.  It is not part of the transform and enters only as the terms $-\ln V-1$ of the
label, \Eq{Sdef}.  In both integrals $V$ is
the large parameter at fixed $P^\mu/V$.  The exponent is stationary where $V\partial\Phi/\partial\beta_\mu=P^\mu$, which is \Eq{targetstate}.
The saddle point is therefore the target state of the particle, $\beta=\bar\beta$.  The value of $V\Phi-\beta_\mu P^\mu$ there is the target entropy
$S_{\rm t}$, and the factor $1/\beta^0$ is the bare temperature $T$ of the target state.  The contour crosses the saddle in the imaginary direction, $\beta=\bar\beta+i\varphi$, where the exponent
falls as $-\tfrac12V\chi^{\mu\nu}\varphi_\mu\varphi_\nu$, and the Gaussian integral over $\varphi$ gives $1/\sqrt{\det(2\pi V\chi)}$.  Hence
\st
\ln G=S_{\rm t}-\tfrac12\ln\det\big(2\pi V\chi^{\mu\nu}\big)+\ln T+O\Big(\frac1{S_{\rm t}}\Big) .
\stp
This is the Gaussian integral of G\"ok\c{c}e's note, used here to go from $Z$ to the entropy instead of from the entropy to $Z$.  The entropy that
appears in the measure is smaller than the thermodynamic entropy by the logarithm of the number of states under the peak.

\paragraph{The scheme.}  We define \emph{state counting at one loop} as the model whose entropy function consists of the first terms of this
expansion.  We keep the two corrections to the target entropy separate, because they have different origins,
\begin{keyeq}
\st
S(e,\bm m,V)=S_{\rm t}-\Delta_\chi-\Delta_\ell,\qquad
\Delta_\chi=\tfrac12\ln\det\big(2\pi V\chi^{\mu\nu}\big),\qquad
\Delta_\ell=\ln\big(e\,\beta^0V\big) ,
\label{scheme}
\stp
\end{keyeq}
where all three terms are evaluated at the target state of $(e,\bm m,V)$, and the $e$ inside the logarithm is the base of the natural logarithm.  The
\emph{determinant term} $\Delta_\chi$ is the logarithm of the number of states under the Gaussian peak.  It comes from the fluctuations of the energy
and the momentum of the particle.  The \emph{label term} $\Delta_\ell$ contains the factor $1/\beta^0$ from the integral over the energy and the factor
$1/(eV)$ that cancels the ideal gas of labels.  We write $\Delta=\Delta_\chi+\Delta_\ell$ for their sum.  Using \Eq{det} for the determinant, the two terms are
\st
\Delta_\chi=\frac D2\ln\big(2\pi V\gamma\,wT_{\rm rel}\big)+\tfrac12\ln a,\qquad
\Delta_\ell=1+\ln V+\ln\gamma-\ln T_{\rm rel} .
\label{Delta}
\stp
The counterterm is evaluated at the target's own state and does not move that state.  This is the difference from the Legendre Hamiltonian, which
puts $\Delta$ inside the transform and thereby shifts the stationary point (\Sect{sec:accuracy}).

\paragraph{The cases.}
\begin{itemize}
\item For any equation of state the two terms are given by \Eq{Delta}, with $w$, $T_{\rm rel}$ and $c_s^2$ taken at the target state.
\item For a constant speed of sound, $p=A\,T_{\rm rel}^{\,n}$ and $q=n-1$, the sum of the two terms can be written with $\mathcal N\equiv1-d\,c_s^2$ as
  \st
  \Delta=\Big(\frac D2+1-\mathcal N\Big)\ln S_{\rm t}+\mathcal N\ln(\gamma V)+\tfrac12\ln a+\text{const},\qquad
  \text{const}=\tfrac D2\ln2\pi-d\,c_s^2\ln(nA)+1 .
  \stp
  The first two terms are the label part and the dynamical part of the split scheme, and $\tfrac12\ln a$ is its sound term.  In the present
  scheme all three are part of the entropy.
\item For the conformal fluid, $q=d$, we have $\mathcal N=0$ and
  \st
  \Delta=\Big(\frac D2+1\Big)\ln S_{\rm t}+\tfrac12\ln a+\text{const} .
  \stp
  Of the coefficient $D/2+1$, the part $D/2$ comes from the determinant term and the $1$ from the label term.  In $d=2$ the coefficient is $5/2$,
  and $\tfrac12\ln a$ is $\tfrac12\ln(1-c_s^2v^2)$ up to a constant.  This is \texttt{SPH\_RENORM=1} with the sound term added as a change of variables.
\end{itemize}

%=============================================================================
\section{The state of a particle}
\label{sec:state}

The temperature, the velocity and the pressure of a particle are the derivatives of $S=S_{\rm t}-\Delta_\chi-\Delta_\ell$ with respect to $e$, $\bm m$ and
$V$, by \Eq{firstlaw}.  Nothing in this section goes beyond differentiating \Eq{scheme} with the chain rule.  The target entropy gives the target's
own thermal vector and pressure.  The two counterterms depend on $e$, $\bm m$ and $V$ only through the target state, so their derivatives are their
slopes in $\beta_\mu$ times the response of the target state.  We first write the result in covariant form, where the response is the inverse
susceptibility, and then in the variables that make the susceptibility diagonal.  In the code the same derivatives could be taken numerically or by
automatic differentiation of $S(e,\bm m,V)$, and that is a check of the closed forms.

\paragraph{The response of the target state.}  The target state is defined by $P^\rho=V\,T^{0\rho}(\bar\beta)$.  Differentiating this condition, with
$\chi^{\rho\mu}=\partial T^{0\rho}/\partial\beta_\mu$, gives
\st
dP^\rho=T^{0\rho}\,dV+V\chi^{\rho\mu}\,d\bar\beta_\mu ,
\stp
so the response of the target state to a change of the energy and momentum, and to a change of the volume, is
\begin{keyeq}
\st
\frac{\partial\bar\beta_\mu}{\partial P^\rho}\bigg|_V=\frac{\chi^{-1}_{\mu\rho}}V,\qquad
\frac{\partial\bar\beta_\mu}{\partial V}\bigg|_P=-\frac{\chi^{-1}_{\mu\rho}\,T^{0\rho}}V .
\label{covresponse}
\stp
\end{keyeq}
The matrix $\chi^{-1}_{\mu\rho}$ is the inverse susceptibility, which converts changes of the conserved densities into changes of $\beta_\mu$ everywhere
in the Density Frame.

\paragraph{The slopes of the two terms.}  Let $\Phi^{\alpha\beta\mu}\equiv\partial\chi^{\alpha\beta}/\partial\beta_\mu$ be the third cumulants of $P^\mu$ per
unit volume.  The slope of the determinant term follows from the derivative of a determinant, and the slope of the label term from
$\beta^0=\eta^{0\mu}\beta_\mu$,
\st
\Delta_\chi^{\,\mu}\equiv\frac{\partial\Delta_\chi}{\partial\beta_\mu}=\tfrac12\,\chi^{-1}_{\alpha\beta}\,\Phi^{\alpha\beta\mu},\qquad
\Delta_\ell^{\,\mu}\equiv\frac{\partial\Delta_\ell}{\partial\beta_\mu}=\frac{\eta^{0\mu}}{\beta^0}=T\,\eta^{0\mu} .
\label{covslopes}
\stp

\paragraph{The thermal vector of the particle.}  By \Eq{firstlaw} the particle has its own thermal vector,
\st
\beta^a_\rho\equiv-\frac{\partial S}{\partial P^\rho}=\frac{(-1,\ \bm v_a)}{T_a} .
\stp
The target entropy contributes $\bar\beta_\rho$, and each counterterm contributes its slope times the response \Eq{covresponse}.  The result is
\begin{keyeq}
\st
\beta^a_\rho=\bar\beta_\rho
+\underbrace{\frac1{2V}\,\chi^{-1}_{\rho\mu}\,\Phi^{\mu\alpha\beta}\,\chi^{-1}_{\alpha\beta}}_{\text{determinant term}}
+\underbrace{\frac TV\,\chi^{-1}_{\rho\mu}\,\eta^{\mu0}}_{\text{label term}} .
\label{betaa}
\stp
\end{keyeq}
The thermal vector of a particle is the target's, shifted by the inverse susceptibility acting on a vector built from the third cumulants and on the
time direction.  Both shifts are of relative order $1/S_{\rm t}$.  The temperature and the velocity of the particle are the components
$1/T_a=-\beta^a_0$ and $\bm v_a/T_a=\bm\beta^a$.

\paragraph{The pressure.}  By \Eq{firstlaw}, $P_a/T_a$ is the derivative of $S$ with respect to the volume at fixed $e$ and $\bm m$.  The target
entropy gives $p\,\bar\beta^0$.  Each counterterm depends on the volume explicitly, through $\tfrac D2\ln V$ and $\ln V$, and through the target state by
\Eq{covresponse}.  The result is
\begin{keyeq}
\st
\frac{P_a}{T_a}=p\,\bar\beta^0
+\underbrace{\frac1V\Big[-\frac D2+\Delta_\chi^{\,\mu}\,\chi^{-1}_{\mu\rho}T^{0\rho}\Big]}_{\text{determinant term}}
+\underbrace{\frac1V\Big[-1+\Delta_\ell^{\,\mu}\,\chi^{-1}_{\mu\rho}T^{0\rho}\Big]}_{\text{label term}} .
\label{covPa}
\stp
\end{keyeq}

\paragraph{The vector $\chi^{-1}_{\mu\rho}T^{0\rho}$ is simple.}  This vector is the change of the target state when the cell is diluted at fixed
energy and momentum.  For a constant speed of sound $\Phi=p\,\beta^0$ is a homogeneous function of $\beta_\mu$ of degree $1-n$, so $T^{0\rho}$ is
homogeneous of degree $-n$.  Euler's theorem then gives $\chi^{\rho\mu}\beta_\mu=-n\,T^{0\rho}$, and therefore
\st
\chi^{-1}_{\mu\rho}\,T^{0\rho}=-\frac{\bar\beta_\mu}n=-\frac pw\,\bar\beta_\mu .
\stp
Diluting the cell rescales $\bar\beta_\mu$ and does nothing else.  The fluid cools and its velocity does not change.  For a general equation of state
there is one more term, which points along the pressure mode defined in \Eq{modes} below.  With $\pi^p_\mu$ the covector of that mode,
\begin{keyeq}
\st
\chi^{-1}_{\mu\rho}\,T^{0\rho}=-\frac pw\,\bar\beta_\mu+\frac{\varepsilon-q\,p}{w\,a}\,\pi^p_\mu .
\label{chiinvT}
\stp
\end{keyeq}
The coefficient $\varepsilon-q\,p=\varepsilon-p/c_s^2$ measures how far the equation of state is from $p=c_s^2\varepsilon$, and it vanishes for a constant
speed of sound.

\paragraph{The modes.}  To make these expressions explicit we use the variables in which the susceptibility is diagonal.  A change of the target state is a change $\delta\tau$ of $\tau\equiv\ln T_{\rm rel}$ and a change $\delta\theta$ of the rapidity.  (A change of the
direction of the velocity does not enter, because $\Delta$ does not depend on the direction.)  The thermal vector changes by
\st
\delta\bar\beta_\mu=-\bar\beta_\mu\,\delta\tau+\frac{\ell_\mu}{T_{\rm rel}}\,\delta\theta,\qquad \ell_\mu=(-\sinh\theta,\ \cosh\theta\,\hat{\bm m}) ,
\stp
and in these variables the susceptibility is
\st
V\chi^{\mu\nu}\,\delta\bar\beta_\mu\,\delta\bar\beta_\nu=S_{\rm t}\,\big(q\,\delta\tau^2+2v\,\delta\tau\,\delta\theta+\delta\theta^2\big) .
\label{quadform}
\stp
Completing the square diagonalizes it,
\st
V\chi^{\mu\nu}\,\delta\bar\beta_\mu\,\delta\bar\beta_\nu=S_{\rm t}\,\big[a\,\delta\tau^2+(\delta\theta+v\,\delta\tau)^2\big],\qquad a=q-v^2 ,
\label{diag}
\stp
with one more square, $\delta u_\perp^2$, for each direction across the velocity.  The energy and momentum change by $\delta P^\mu=V\chi^{\mu\nu}\delta\bar\beta_\nu$.  The combinations of $\delta e$ and
$\delta\bm m$ conjugate to $\delta\tau$, to $\delta\theta+v\,\delta\tau$ and to $\delta u_\perp$ are, in units of $T_{\rm rel}$,
\begin{keyeq}
\st
\phi_p=\frac{\gamma\big[(1+v^2)\,\delta e-2v\,\delta m_\parallel\big]}{T_{\rm rel}},\qquad
\phi_\parallel=\frac{\gamma\,(\delta m_\parallel-v\,\delta e)}{T_{\rm rel}},\qquad
\phi_{\perp\alpha}=\frac{\delta m_{\perp\alpha}}{T_{\rm rel}} ,
\label{modes}
\stp
\st
\phi_p=S_{\rm t}\,a\,\delta\tau,\qquad \phi_\parallel=S_{\rm t}\,(\delta\theta+v\,\delta\tau),\qquad
\avg{\phi_p^2}=S_{\rm t}\,a,\qquad \avg{\phi_\parallel^2}=\avg{\phi_{\perp\alpha}^2}=S_{\rm t} .
\label{response}
\stp
\end{keyeq}
They are uncorrelated, and each responds to its own variable only.  We write $\pi^p_\mu$ and $\pi^\parallel_\mu$ for the covectors of the first two
modes, so that $\phi_p=\pi^p_\mu\,\delta P^\mu$ and $\phi_\parallel=\pi^\parallel_\mu\,\delta P^\mu$,
\st
\pi^p_\mu=\frac{\gamma\,\big(1+v^2,\ -2v\,\hat{\bm m}\big)}{T_{\rm rel}},\qquad
\pi^\parallel_\mu=\frac{\gamma\,\big(-v,\ \hat{\bm m}\big)}{T_{\rm rel}} .
\stp
In terms of them the inverse susceptibility is a sum of one term per mode,
\begin{keyeq}
\st
\frac{\chi^{-1}_{\rho\mu}}V=\frac1{S_{\rm t}}\Big[\frac{\pi^p_\rho\,\pi^p_\mu}a+\pi^\parallel_\rho\,\pi^\parallel_\mu+\dots\Big] ,
\label{chiinv}
\stp
\end{keyeq}
where the dots stand for the transverse directions, which have the same form as the second term.
\begin{itemize}
\item $\phi_p$ is the \emph{pressure mode}.  Since $\delta p=w\,\delta\tau$, it measures the change of the pressure, which is the Density Frame
  $\delta p$ at given $\delta e$ and $\delta\bm m$.  Its stiffness is the acoustic norm,
  \st
  \delta p=\frac{(1+v^2)\,\delta e-2v\,\delta m_\parallel}{V\,a},\qquad \avg{(\delta\ln T_{\rm rel})^2}=\frac1{S_{\rm t}\,a} .
  \stp
\item $\phi_\parallel$ and $\phi_\perp$ are the \emph{comoving momenta}.  They are the momentum of the change as seen in the rest frame of the target, along and
  across the velocity.
\end{itemize}
The determinant is the product of the variances.  The change of variables from $\delta P^\mu$ to the modes has the Jacobian $T_{\rm rel}^{-D}$, so the determinant term of \Eq{Delta} is that of $D$ Gaussian modes,
\st
\det(V\chi)=T_{\rm rel}^{2D}\,S_{\rm t}^{D}\,a,\qquad
\Delta_\chi=\frac D2\ln\big(2\pi\,S_{\rm t}T_{\rm rel}^2\big)+\tfrac12\ln a .
\stp

\paragraph{The shifts.}  In the modes \Eq{betaa} becomes one division per mode.  Contracting the slopes \Eq{covslopes} with the covectors of the modes
defines the slopes of the counterterm along the modes,
\st
\Delta_p\equiv\pi^p_\mu\,\Delta^\mu=\frac{\partial\Delta}{\partial\tau}-v\,\frac{\partial\Delta}{\partial\theta},\qquad
\Delta_\theta\equiv\pi^\parallel_\mu\,\Delta^\mu=\frac{\partial\Delta}{\partial\theta} .
\label{slopes}
\stp
The first is the slope of the counterterm in $\tau$ when $\theta+v\tau$ is held fixed, and the second is its slope in $\theta+v\tau$ when $\tau$ is held
fixed.  These are the natural slopes because $\phi_p$ determines $\delta\tau$ and $\phi_\parallel$ determines $\delta\theta+v\,\delta\tau$.  With \Eq{chiinv} the
thermal vector of the particle is
\begin{keyeq}
\st
\beta^a_\rho=\bar\beta_\rho+\epsilon_p\,\pi^p_\rho+\epsilon_\parallel\,\pi^\parallel_\rho,\qquad
\epsilon_p=\frac{\Delta_p}{S_{\rm t}\,a},\qquad \epsilon_\parallel=\frac{\Delta_\theta}{S_{\rm t}} .
\label{shift}
\stp
\end{keyeq}
Each shift is the slope of the counterterm along a mode divided by the variance of that mode.  Equivalently, the entropy of the particle changes at
fixed volume by
\st
dS=-\beta^a_\rho\,dP^\rho=(1-\epsilon_p)\,\phi_p+(v-\epsilon_\parallel)\,\phi_\parallel .
\stp
In the language of thermodynamics the shifts are the entropic conjugates of the modes in the counterterm, just as $1$ and $v$ are their conjugates
in the target entropy,
\st
d\Delta=\epsilon_p\,\phi_p+\epsilon_\parallel\,\phi_\parallel,\qquad dS_{\rm t}=\phi_p+v\,\phi_\parallel .
\stp
The components of \Eq{shift} give the temperature and the velocity,
\begin{keyeq}
\st
\frac1{T_a}=\frac{\gamma}{T_{\rm rel}}\big[1-(1+v^2)\,\epsilon_p+v\,\epsilon_\parallel\big],\qquad
v_a=\frac{v\,(1-2\epsilon_p)+\epsilon_\parallel}{1-(1+v^2)\,\epsilon_p+v\,\epsilon_\parallel},\qquad \bm v_a\parallel\bm m .
\label{Tava}
\stp
\end{keyeq}
To first order in the shifts this is the target's thermal vector at the temperature $T_{\rm rel}(1+\epsilon_p)$ and the rapidity
$\theta+\epsilon_\parallel-v\,\epsilon_p$.  \Eq{Tava} itself is exact for the scheme \Eq{scheme}.

\paragraph{The slopes of the two terms in the modes.}  To find the slopes we write \Eq{Delta} as functions of $\tau$, $\theta$ and $V$.  With
$\gamma=\cosh\theta$, $v=\tanh\theta$ and $T_{\rm rel}=e^{\tau}$ the two terms are
\st
\Delta_\chi=\frac D2\Big[\ln\big(2\pi V\,wT_{\rm rel}\big)+\ln\cosh\theta\Big]+\tfrac12\ln\big(q-\tanh^2\theta\big),\qquad
\Delta_\ell=1+\ln V+\ln\cosh\theta-\tau .
\stp
The equation of state enters only the determinant term, through $wT_{\rm rel}$ and $q$.  They depend on $\tau$ alone, with
$d\ln(wT_{\rm rel})/d\tau=2+q$ and $dq/d\tau=\dot q$.  Differentiating each term and forming the combination \Eq{slopes} gives
\begin{keyeq}
\begin{align}
\Delta_p^{\chi}&=\frac D2\,(2+q-v^2)+\frac{\dot q+2v^2(1-v^2)}{2a}, & \Delta_p^{\ell}&=-(1+v^2), \nonumber\\
\Delta_\theta^{\chi}&=\frac D2\,v-\frac{v\,(1-v^2)}a, & \Delta_\theta^{\ell}&=v .
\label{Dderiv}
\end{align}
\end{keyeq}
The slopes in \Eq{shift} are the sums $\Delta_p=\Delta_p^\chi+\Delta_p^\ell$ and $\Delta_\theta=\Delta_\theta^\chi+\Delta_\theta^\ell$, so the shifts
separate as
\st
\epsilon_p=\frac{\Delta_p^{\chi}}{S_{\rm t}\,a}-\frac{1+v^2}{S_{\rm t}\,a},\qquad
\epsilon_\parallel=\frac{\Delta_\theta^{\chi}}{S_{\rm t}}+\frac v{S_{\rm t}} .
\stp
The two terms pull in opposite directions on the temperature.  The determinant term makes a particle hotter than its target state, and the label
term makes it cooler.  The label term contains no derivative of the equation of state.  At rest $\Delta_\theta=0$ and
$\Delta_p=\tfrac D2(2+q)+\dot q/2q-1$, where the $-1$ is the label term, so a particle at rest is shifted in temperature only.

\paragraph{The pressure in the modes.}  With \Eq{chiinvT} the contractions in \Eq{covPa} are $\Delta^\mu\bar\beta_\mu=-\partial\Delta/\partial\tau$ and
$\Delta^\mu\pi^p_\mu=\Delta_p$.  The pressure of the particle is then
\begin{keyeq}
\st
\frac{P_aV}{T_a}=\frac{p\,V\gamma}{T_{\rm rel}}
+\underbrace{\Big[-\frac D2+\frac pw\,\frac{\partial\Delta_\chi}{\partial\tau}+\frac{\varepsilon-q\,p}{w\,a}\,\Delta_p^{\chi}\Big]}_{\text{determinant term}}
+\underbrace{\Big[-1-\frac pw-\frac{\varepsilon-q\,p}{w\,a}\,(1+v^2)\Big]}_{\text{label term}} ,
\label{Pa}
\stp
\end{keyeq}
where $\partial\Delta_\chi/\partial\tau=\tfrac D2(2+q)+\dot q/2a$.  For a constant speed of sound the terms with $\varepsilon-q\,p$ vanish and $p/w=1/n$.

\paragraph{Why the pressure of the model is the target's.}  The volume enters the single-particle integral \Eq{condition} only through $H$, so the
derivative of $\ln Z_1$ with respect to the volume is the mean of the particle's pressure.  With the right side of \Eq{condition} for $Z_1$ it is
\st
\avg{P_a}=T\,\frac{\partial\ln Z_1}{\partial V}=p-\frac TV .
\stp
The pressure of the model is the mean particle pressure plus the pressure $\bar nT$ of the ideal gas of
labels (\nolinkurl{partition.tex}).  At the mean volume, $V=1/\bar n$, the term $-T/V$ cancels the pressure of the labels, and what is left is $p$.  In
the one-loop scheme $Z_1$ differs from \Eq{condition} by the factor $e^{-c/S_{\rm t}}$ of \Sect{sec:accuracy}.  The cancellation of the label pressure
is unchanged, and the pressure is off by the relative amount $n\,c/S_{\rm t}^2$.

\paragraph{The Hamiltonian.}  Given the label, the momentum and the volume of a particle, its energy $H(s,\bm m,V)$ is the solution $e$ of
\st
S(e,\bm m,V)=s .
\label{root}
\stp
This is one equation in one unknown.  Its derivative $\partial S/\partial e=1/T_a$ is known from \Eq{Tava}, so Newton's iteration is
\st
e\;\to\;e+T_a\,\big[s-S(e,\bm m,V)\big] .
\label{newton}
\stp

\paragraph{The domain.}  $H$ exists where $S$ increases with $e$ and $|\bm v_a|<1$, that is, where the thermal vector $\beta^a_\mu$ of the particle is
timelike.  The condition is
\st
1-\epsilon_p>|\epsilon_\parallel-v\,\epsilon_p| .
\stp
At rest this condition is $S_{\rm t}>\Delta_p/q$, which is $S_{\rm t}>D/2+1$ for the conformal fluid; this is the label boundary of
\texttt{SPH\_RENORM=1}.  Below it the saddle point has no meaning, and a state there has no Hamiltonian, as in the present schemes.  The boundary
belongs to the truncation of the series, because the exact $G$ increases with $e$ by construction.

Equations \Eq{Tava} and \Eq{Pa} were compared with finite differences of \Eq{scheme} for a two-component $p=T^3+0.7\,T^5$
(\nolinkurl{examples/renorm_count_state.py}).

%=============================================================================
\section{In the code}
\label{sec:code}

The variables $(\bm x_a,\bm m_a,s_a)$, the Liouville measure and both steps of the scheme stay as they are.  The new scheme enters only through the
two particle-level maps of \nolinkurl{src/pointwise.jl}, in the same way as \nolinkurl{renorm.jl} and \nolinkurl{renorm_split.jl} do.  Everything
a step needs follows once the target state $\bar\beta_\mu$ of the particle is known, so the two maps differ only in how they find it.

\paragraph{From the energy.}  When the energy is known, as in the viscous step, nothing has to be solved beyond what the code does now.  Given $e$,
$\bm m$ and $V$, we first find $T_{\rm rel}$ and $v$ of the target state with the standard equation of state.  We can then evaluate the label
$s=S_{\rm t}-\Delta_\chi-\Delta_\ell$ from \Eq{Delta}, and finally the thermal vector of the particle from \Eq{betaa}, or equivalently its temperature and
velocity from \Eq{Tava}.  For a state outside the domain of \Sect{sec:state} the map returns $s=-\infty$.

\paragraph{From the label.}  In the ideal step the label is known and the energy is not.  The target state is defined by the energy and the
momentum, so it cannot be found before the energy.  Finding the target state from $(s,\bm m,V)$ and solving \Eq{root} for the energy are the same
problem, and there are two equivalent ways to solve it.
\begin{itemize}
\item We can iterate on the energy alone with Newton's step \Eq{newton}, starting from the previous energy of the particle.  Each iteration calls the
  map from the energy once, and it reduces the error by a factor of order $1/S_{\rm t}$.
\item We can instead solve directly for $\tau$ and $\theta$ of the target state from the two equations
  \st
  s=\gamma\,s(T_{\rm rel})\,V-\Delta_\chi(\tau,\theta,V)-\Delta_\ell(\tau,\theta,V),\qquad |\bm m|=V\,w\,\gamma^2v ,
  \stp
  by a Newton iteration in two unknowns, as \nolinkurl{renorm.jl} does for the Legendre Hamiltonian.  No inner inversion is then needed.
\end{itemize}
In either case the energy is a by-product, $e=V(w\gamma^2-p)$, and requires no further iteration.  Once the target state is known, the thermal vector
$\beta^a_\rho$ and the pressure $P_a$ follow directly from \Eq{betaa} and \Eq{Pa}.

\paragraph{The ideal step} is the same generalized leapfrog at fixed labels.  For each particle we find the target state from $(s_a,\bm m_a,V_a)$ as
just described, and then use the velocity $\dot{\bm x}_a=\bm v_a$ and build the pressure force from $P_aV_a^2$ in place of $p_aV_a^2$.  The leapfrog is
implicit, so the solve is repeated at each of its inner iterations, starting each time from the previous target state.  The Hamiltonian of the step
is $H=\sum_aH(s_a,\bm m_a,V_a)$, so it conserves $H$ and $\bm P$ and preserves the measure $\prod_ad\bm x_a\,d\bm m_a$.  It therefore keeps the
distribution
\st
\pi\propto\exp\sum_a\Big[s_a-\frac{H_a-\bm v\cdot\bm m_a}T\Big] .
\stp

\paragraph{The viscous step} moves the momenta at fixed positions and energies.  The energy $e_a$ of each particle is known from the last solve of
the ideal step, so nothing has to be found before the step begins.  Its target is the same distribution written in the energies,
\st
\pi\propto\exp\sum_a\Big[S(e_a,\bm m_a,V_a)-\ln T_a\Big] .
\stp
For a proposed transfer of momentum $\Delta\bm m$ we evaluate $S(e_a,\bm m_a+\Delta\bm m,V_a)-\ln T_a$ with the map from the energy at the new
momentum.  The proposals, the Hastings factor and the sweeps do not change.  At the end of the step the label of each particle is set to
$S(e_a,\bm m_a,V_a)$ with its final momentum.

\paragraph{One function for both steps.}  The weight in the energies must be
\st
\ln\frac{\partial e^{S}}{\partial e}=S-\ln T_a ,
\stp
with the temperature $T_a$ derived from the same function $S$.  It should not be a separate one-loop formula for $\ln g$.  With this choice both steps keep exactly the same
distribution, whatever the order to which $S$ approximates \Eq{Sdef}.  The truncation of the series then affects the thermodynamics of the model
but not its detailed balance.

\paragraph{The rest of the code.}  Hybrid Monte Carlo uses the same Hamiltonian.  The measured energy is $\sum_aH_a$, the virial uses the particle
pressure $P_a$, and the kinetic stress is $\bm v_a\bm m_a$.  The equation of state must supply $p$, $s$, $c_s^2$ and $\dot q$ at $T_{\rm rel}$.  The
transport coefficients are evaluated as they are now.

%=============================================================================
\section{The accuracy of one loop}
\label{sec:accuracy}

The label of the scheme differs from the exact one by the next term of the series (\App{sec:next}),
\st
\ln G-\ln V-1-S=\frac{c}{S_{\rm t}}+O\Big(\frac1{S_{\rm t}^2}\Big),
\stp
so the single-particle partition function of the model is slightly off.  For a constant speed of sound the result is the same relative error in
the energy and the pressure,
\st
\ln Z_1=\ln Z_{\rm t}-\ln(eV)-\frac c{S_{\rm t}},\qquad \frac{\delta E}E=\frac{\delta p}p=\frac{n\,c}{S_{\rm t}^2} .
\stp  \Tab{tab:acc} shows the quadrature of the single particle at rest in $d=2$
(\nolinkurl{examples/renorm_count_quadrature.py}).

The Legendre Hamiltonian instead takes the extremum of $V\Phi-\Delta-\beta_\mu P^\mu$.  Because $\Delta$ depends on $\beta$, this moves the
stationary point, and the label is lowered by
\st
\frac1{2V}\,\Delta^\mu\,\chi^{-1}_{\mu\nu}\,\Delta^\nu,\qquad \Delta^\mu\equiv\frac{\partial\Delta}{\partial\beta_\mu} .
\stp
For the determinant part of $\Delta$ this is the second term of $C_g$ in \Eq{CgCh}, taken without the other two, which nearly cancel it.  In place of $c$ the Legendre scheme
therefore has the coefficient
\st
\kappa=\frac{(3n+1)^2}{8(n-1)}+c ,
\stp
and $n\kappa/S_{\rm t}^2$ reproduces its measured residuals
(\nolinkurl{renorm2.tex}).

\begin{table}[h]
\caption{The relative error of the pressure and the energy of one particle at rest in $d=2$.  The columns are state counting at one loop from the
quadrature, its leading term $n\,c/S_{\rm t}^2$, and the Legendre Hamiltonian, $n\kappa/S_{\rm t}^2$.}
\label{tab:acc}
\begin{ruledtabular}
\begin{tabular}{crccc}
$c_s^2$ & $S_{\rm t}$ & one loop, quadrature & $n\,c/S_{\rm t}^2$ & Legendre \\
\hline
$1/2$ & 54  & $1.7\times10^{-4}$ & $0.9\times10^{-4}$ & $6.5\times10^{-3}$ \\
      & 541 & $9.2\times10^{-7}$ & $8.5\times10^{-7}$ & $6.5\times10^{-5}$ \\
$1/4$ & 54  & $5.1\times10^{-4}$ & $3.2\times10^{-4}$ & $1.4\times10^{-2}$ \\
      & 541 & $3.3\times10^{-6}$ & $3.2\times10^{-6}$ & $1.4\times10^{-4}$ \\
\end{tabular}
\end{ruledtabular}
\end{table}
Two checks have not been made yet.  The one-loop scheme has not been tested in motion, and it has not been compared with the quadrature of the split scheme.

%=============================================================================
\appendix

\section{The next order of the expansion}
\label{sec:next}

This appendix records the term of order $1/V$.  It is not needed for state counting at one loop, and we do not plan to implement it soon.


\paragraph{The expansion.}  Along the contour the exponent is
\st
S_{\rm t}-\frac V2\,\chi^{\mu\nu}\varphi_\mu\varphi_\nu-\frac{iV}{3!}\,\Phi^{\mu\nu\rho}\varphi_\mu\varphi_\nu\varphi_\rho+\frac V{4!}\,\Phi^{\mu\nu\rho\sigma}\varphi_\mu\varphi_\nu\varphi_\rho\varphi_\sigma+\dots,
\stp
where $\Phi^{\mu\cdots}$ are the derivatives of $\Phi$ with respect to $\beta_\mu$.  They are the cumulants of $P^\mu$ per unit volume.  The term of
order $1/V$ comes from the quartic vertex taken once and the cubic vertex taken twice.  For the label there is in addition the factor $1/\beta^0$,
which we write as $e^{h}$ with $h=-\ln\beta^0$ and $h^\mu=\partial h/\partial\beta_\mu$.  The result is
\st
\ln g=S_{\rm t}-\tfrac12\ln\det(2\pi V\chi)+\frac{C_g}V,\qquad \ln G=\ln g+\ln T+\frac{C_h}V ,
\stp
\begin{align}
C_g&=\tfrac18\,\Phi^{\mu\nu\rho\sigma}\chi^{-1}_{\mu\nu}\chi^{-1}_{\rho\sigma}
-\tfrac18\,\Phi^{\mu\nu\rho}\Phi^{\alpha\beta\gamma}\chi^{-1}_{\mu\nu}\chi^{-1}_{\rho\alpha}\chi^{-1}_{\beta\gamma}
-\tfrac1{12}\,\Phi^{\mu\nu\rho}\Phi^{\alpha\beta\gamma}\chi^{-1}_{\mu\alpha}\chi^{-1}_{\nu\beta}\chi^{-1}_{\rho\gamma}, \nonumber\\
C_h&=-\tfrac12\big(h^{\mu\nu}+h^\mu h^\nu\big)\chi^{-1}_{\mu\nu}+\tfrac12\,h^\alpha\chi^{-1}_{\alpha\mu}\Phi^{\mu\nu\rho}\chi^{-1}_{\nu\rho} .
\label{CgCh}
\end{align}
Each derivative of $\Phi$ is of order $s\,T_{\rm rel}^{\,k}$, so $C/V=c/S_{\rm t}$ with $c$ of order one.  The series in $1/V$ is thus a series in
$1/S_{\rm t}$.  All the derivatives follow from one tower of functions of the temperature,
\st
b_1=p,\qquad b_{k+1}=T^3\frac{db_k}{dT},\qquad \frac{\partial^k\Psi}{\partial\beta_{\mu_1}\cdots\partial\beta_{\mu_k}}=\sum_jb_{k-j}\,\big[g^{\,j}\beta^{\,k-2j}\big]^{\mu_1\cdots\mu_k}_{\rm sym} ,
\stp
where the bracket is the sum over the distinct pairings.  The order $1/V$ needs $b_5$,
which means that it needs $q$, $\dot q$ and $\ddot q$.

\paragraph{The variables.}  The cubic and quartic terms are simplest in the modes of \Eq{modes}, the pressure mode $\phi_p$ and the comoving momenta
$\phi_\parallel$, $\phi_{\perp\alpha}$ ($\alpha=1,\dots,d-1$), which diagonalize the quadratic term.

\paragraph{The cumulants.}  Let $k_{ij\cdots}$ be the joint cumulants of $(\phi_p,\phi_\parallel,\phi_\perp)$ divided by $S_{\rm t}$, with the indices
$p$, $\parallel$ and $\perp$.  Each of them is the corresponding cumulant in the rest frame plus a polynomial in $v^2$, and they are the same in any
dimension.  We write the rest-frame cumulants of the energy, and of $\delta e^2m_i^2$, as
\st
\hat\varepsilon_3=q^2+2q+\dot q,\qquad \hat\varepsilon_4=q^3+5q^2+6q+(3q+5)\dot q+\ddot q,\qquad \hat m_4=(q+2)(q+3)+\dot q ,
\stp
\begin{align}
k_{pp}&=a, & k_{\parallel\parallel}&=k_{\perp\perp}=1, \nonumber\\
k_{ppp}&=\hat\varepsilon_3-3v^4, & k_{p\parallel\parallel}&=q+2-3v^2, & k_{pp\parallel}&=-v\,k_{p\parallel\parallel}, \nonumber\\
k_{p\perp\perp}&=q+2-v^2, & k_{\parallel\parallel\parallel}&=3v, & k_{\parallel\perp\perp}&=v, \nonumber\\
k_{pp\parallel\parallel}&=\hat m_4-3v^2(q+3), & k_{pp\perp\perp}&=\hat m_4-v^2(q+3), \nonumber\\
k_{pppp}&=\hat\varepsilon_4+2v^2\hat m_4-9v^4(q+3), & k_{\parallel\parallel\parallel\parallel}&=k_{\perp\perp\perp\perp}=3(q+3), &
k_{\parallel\parallel\perp\perp}&=k_{\perp_1\perp_1\perp_2\perp_2}=q+3 .
\end{align}

\paragraph{The order $1/V$ in these variables.}  We define $c_g=C_gS_{\rm t}/V$, $c_h=C_hS_{\rm t}/V$ and $w_i\equiv\sum_jk_{ijj}/k_j$.  In the modes
the susceptibility is diagonal, so \Eq{CgCh} no longer contains an inverse matrix,
\begin{keyeq}
\st
c_g=\frac18\sum_{i,j}\frac{k_{iijj}}{k_ik_j}-\frac18\sum_i\frac{w_i^2}{k_i}-\frac1{12}\sum_{i,j,k}\frac{k_{ijk}^2}{k_ik_jk_k},\qquad
c_h=-\frac{(1+v^2)^2}a-v^2+\frac12\Big[\frac{(1+v^2)\,w_p}a-v\,w_\parallel\Big] ,
\label{cdiag}
\stp
\end{keyeq}
where the sums run over $p$, $\parallel$ and the $d-1$ transverse directions.  The three combinations $w_i$ are
\st
w_p=\frac{k_{ppp}}a+k_{p\parallel\parallel}+(d-1)\,k_{p\perp\perp},\qquad w_\parallel=\frac{k_{pp\parallel}}a+(d+2)\,v,\qquad w_\perp=0 .
\stp

\paragraph{The closed form.}  When the result is written in terms of $a$, the velocity enters only through powers of $1/a$.  For the label the
coefficient is $c=c_g+c_h$, with
\begin{keyeq}
\st
c=-\frac{d^2-1}8-\frac{(d-2)(q+d-1)}{2a}+\frac{\lambda_4}{8a^2}-\frac{5\lambda_3^2}{24a^3},
\label{c}
\stp
\st
\lambda_3=2q(q-1)-\dot q,\qquad \lambda_4=2q(q-1)(3q+4d-7)-(7q+4d-9)\,\dot q+\ddot q .
\stp
\end{keyeq}
The last two terms have the form of the correction for a single variable with the cumulants $\kappa_2=a$, $\kappa_3=\lambda_3$ and
$\kappa_4=\lambda_4$, which is
\st
\frac{\kappa_4}{8\kappa_2^2}-\frac{5\kappa_3^2}{24\kappa_2^3} .
\stp  After the momenta are integrated, what remains is the pressure mode.  The two parts are separately
\begin{align}
c_h&=-\frac{d-1}2+\frac{2(q+1)(q+d-1)-\dot q}{2a}-\frac{(q+1)\lambda_3}{2a^2}, \nonumber\\
c_g&=-\frac{(d-1)(d-3)}8-\frac{d(d+3q-1)+\lambda_3}{2a} \nonumber\\
&\quad+\frac{2q(q-1)(7q+4d-3)-(11q+4d-5)\,\dot q+\ddot q}{8a^2}-\frac{5\lambda_3^2}{24a^3} .
\end{align}
The cubic coefficient $\lambda_3$ is the part of $-k_{ppp}$ that does not vanish with $a$.  In terms of the speed of sound it is
\st
\frac{\lambda_3}{2q^2}=(1-c_s^2)+\frac w{c_s}\frac{dc_s}{de} ,
\stp
which appears to be the steepening coefficient of a relativistic simple wave.  We have not checked this
identification against the literature.

\paragraph{The cases.}
\begin{itemize}
\item For a constant speed of sound, $q=n-1$ and $\dot q=\ddot q=0$, so that $\lambda_3=2q(q-1)$ and $\lambda_4=2q(q-1)(3q+4d-7)$.
\item In $d=2$ the term in $1/a$ vanishes, and
  \st
  c=-\frac38+\frac{\lambda_4}{8a^2}-\frac{5\lambda_3^2}{24a^3} .
  \stp
  At rest this is $c=1/12$ for $c_s^2=1/2$ and $c=3/16$ for $c_s^2=1/4$.
\item For the conformal fluid, $q=d$, the coefficients are $\lambda_3=2d(d-1)$ and $\lambda_4=14\,d\,(d-1)^2$, with $a=d-v^2$.
\end{itemize}

\paragraph{State counting at two loops} would be \Eq{scheme} with the two terms corrected separately,
\st
\Delta_\chi\;\to\;\Delta_\chi-\frac{c_g}{S_{\rm t}},\qquad \Delta_\ell\;\to\;\Delta_\ell-\frac{c_h}{S_{\rm t}} ,
\stp
since $c_g$ belongs to the determinant and $c_h$ to the factor $1/\beta^0$.  Everything in \Sect{sec:state} and \Sect{sec:code} then holds with the new
terms in the slopes.  Their derivatives follow from
\st
\frac{\partial\ln S_{\rm t}}{\partial(\tau,\theta,\ln V)}=(q,\,v,\,1),\qquad \frac{\partial a}{\partial(\tau,\theta)}=\big(\dot q,\,-2v(1-v^2)\big) ,
\stp
and they need the third derivative of $q$.  The remaining error in the label would be of order $1/S_{\rm t}^2$.

\paragraph{The checks.}  \Eq{c} is an identity with \Eq{cdiag}, for any $d$.  We compared the tensor contraction \Eq{CgCh} with it in three ways.  For
a general equation of state the two agree symbolically for $d=1,2,3$ at any $v$.  For $p=T^3+0.7\,T^5$ they agree with the direct derivatives of
$p\,\beta^0$.  For a constant speed of sound in $d=2$ the exact series for $G$ gives $(\ln G-\text{one loop})\,S_{\rm t}$, which converges to $c(v)$
to three digits at $S_{\rm t}=541$ (\nolinkurl{examples/renorm_saddle_check.py}, \nolinkurl{renorm2.tex}).

\end{document}
